\begin{tabular}{l l rr r p{10.4cm}}
\toprule
Model & Run & val MASE & val sMAPE & s & Rationale for the run \\
\midrule
SARIMAX & s1 & 0.717 & 12.1 & 3 & Minimal seasonal random walk + AR(1): the simplest structure consistent with the ACF (slow decay at multiples of 24). Reference point for what the seasonal difference alone buys. \\
SARIMAX & s2 & 0.659 & 10.8 & 44 & Add MA terms at lag 1 and 24: PACF after seasonal differencing shows a negative spike at 24, the classic signature of a seasonal MA(1). \\
SARIMAX & s3 & 0.666 & 11.1 & 108 & Weekly cycle cannot be modelled with a 168-period SARIMA (state too large). Add 3 Fourier pairs of period 168 h as exogenous regressors (dynamic harmonic regression). \\
SARIMAX & s4 & 0.669 & 11.2 & 139 & Public holidays behave like Sundays; a holiday dummy lets the level drop without disturbing the seasonal state. \\
SARIMAX & s5 & 0.769 & 12.7 & 416 & Check whether more harmonics (5) and a second AR lag improve or overfit; AIC vs validation MASE decides. \\
LightGBM & g1 & 0.642 & 10.7 & 7 & Lags 1-24 only: how much can trees recover from the last day alone? Establishes the value of the seasonal/weekly lags added next. \\
LightGBM & g2 & 0.672 & 10.8 & 7 & Add lags 36-168 (two-day and weekly cycle) identified in the ACF. \\
LightGBM & g3 & 0.604 & 9.7 & 10 & Add calendar features of origin and target (hour, weekday, holiday) so the model can anticipate the day type of the target time, which lags alone cannot at 24 h. \\
LightGBM & g4 & 0.664 & 10.6 & 9 & Add rolling means/stds (6, 24, 168 h) and differences: level and volatility summaries that reduce reliance on any single lag. \\
LightGBM & g5 & 0.682 & 10.9 & 6 & Regularise: fewer leaves, lower learning rate and larger leaves to test whether g4 overfits the pooled data (compare train/val gap). \\
LightGBM & g6 & 0.652 & 10.4 & 9 & Opposite direction: more capacity. Together with g5 this brackets the capacity sweet spot before the Optuna search. \\
LightGBM & Optuna (20 trials) & 0.633 & 10.2 & 244 & best trial opt05: num\_leaves=67, learning\_rate=0.0182, min\_child\_samples=32, feature\_fraction=0.796, lambda\_l2=0.00153 \\
RNN & l1 & 0.721 & 11.5 & 12 & Small network with one day of context, to see the floor of the architecture and whether training is stable. \\
RNN & l2 & 0.747 & 12.5 & 6 & One week of context: the weekly cycle should now be visible to the network. \\
RNN & l3 & 0.714 & 11.5 & 6 & Double the hidden size to test whether l2 was capacity-limited. \\
RNN & l4 & 0.793 & 12.6 & 7 & Deeper network with dropout: more capacity, but regularised against the small number of training weeks. \\
RNN & l5 & 0.711 & 11.6 & 7 & GRU has fewer parameters than LSTM at equal hidden size; test whether the gating simplification helps with ~6 weeks of data. \\
RNN & l6 & 0.720 & 11.4 & 12 & Lower learning rate and light weight decay to smooth the validation curve of the best previous configuration. \\
\bottomrule
\end{tabular}
